\documentclass[10pt,a4paper]{amsart}
\usepackage[margin=19mm]{geometry}
\usepackage{amsmath,amssymb,mathtools}
\usepackage[scr=rsfs,cal=euler]{mathalfa}
\usepackage{xfrac}
\usepackage[unicode,hidelinks,bookmarks=false]{hyperref}
\newcommand{\ie}{i.e.\ }
\newcommand{\cf}{cf.\ }
\newcommand{\resp}{respectively\ }
\newcommand{\Subsection}{\subsection}
\providecommand{\nobreakdash}{\nobreak\hbox{-}}
\setlength{\emergencystretch}{2em}
\multlinegap0pt
\usepackage{amssymb}
\newtheorem{lemma}{Lemma}
\def\RRR{\mathfrak{R}}
\def\gl{\mathfrak{gl}}
\def\sl{\mathfrak{sl}}
\def\zzz{\mathfrak{z}}
\title{Uniqueness of addition in Lie rings $\gl_n(K)$ and $\sl_n(K)$}
\author{Gennadiy Sosnov}
\date{Source: arXiv:2606.08201v1 (CC BY 4.0)}
\begin{document}
\maketitle

\noindent\emph{Source and licence.} Uniqueness of addition in Lie rings $\gl_n(K)$ and $\sl_n(K)$. Version arXiv:2606.08201v1 (CC BY 4.0), distributed by arXiv under the Creative Commons Attribution 4.0 International licence, http://creativecommons.org/licenses/by/4.0/, as recorded in the arXiv licence metadata for this exact version and shown on the abstract page https://arxiv.org/abs/2606.08201v1. Full licence notice: \texttt{licenses/arxiv-2606.08201-CC-BY-4.0.txt}.\par\smallskip

\noindent\textit{Editorial scope.} Counted unit: the article's lemma lemma:almostAdditivityForDiagInSl (source0.tex lines 826--828). The article's complete original proof of it is delivered with it (source0.tex lines 830--859, one \texttt{proof} environment of 30 lines). No mathematics is written, repaired or re-derived: the statement and the proof are the article's own lines, in the article's own notation, and no retained line is substituted or re-flowed.  Retained uncounted as the source-local context the proof needs: lines 214--219; lines 335--337.  Nothing was omitted from any retained span, so the package carries no omission record.  No retained line contains a citation, so the wrapper carries no bibliography and no bibliography entry.  No retained line contains a figure environment, an image-inclusion command or drawing code, so no image is delivered and the package declares no asset dependency.  Editorial formatting only, in addition: an AMS \texttt{amsart} preamble that restates the article's own declarations and macros used by the retained lines, namely the environments \texttt{lemma} on separate counters, together with the standard packages those lines need.  The retained environments print in this document as Lemma 1 in order of appearance, whereas the article prints the same items with the numbering of its own full document.  The complete native source of the article as distributed in the arXiv e-print for this exact version is bundled unchanged as \texttt{sources/202315/source0.tex}.
	\begin{lemma}\label{lemma:intersectionOfCenters}
		For any $a, b \in \RRR$ we have
		\begin{equation*}
			\Delta(a, b) \in \alpha(Z(C(a)) \cap Z(C(b))).
		\end{equation*}
	\end{lemma}

	\begin{lemma}\label{lemma:additivityForBothDiag}
		Let $A, B \in \gl_n(K)$ be two diagonal matrices. Then $\Delta(A, B) \in \alpha(\zzz)$.
	\end{lemma}

	\begin{lemma}\label{lemma:almostAdditivityForDiagInSl}
		Let $A, B \in \sl_n(K)$ and $B$ be a diagonal matrix. Then $\Delta(A, B) \in \alpha(\zzz)$.
	\end{lemma}

	\begin{proof}
		Again, it is sufficient to prove the statement in the case $B = \lambda (E_{ii} - E_{jj})$. If $A$ is a diagonal matrix, then the lemma follows from Lemma~\ref{lemma:additivityForBothDiag}. Therefore, without loss of generality, we can assume that $A$ is not a diagonal matrix. Let $A_{st} \neq 0$, where $s \neq t$. If $Z(C(A)) \cap Z(C(B)) = \zzz$, then the lemma follows from Lemma~\ref{lemma:intersectionOfCenters}. Otherwise, $Z(C(B)) \subseteq Z(C(A))$, so $[A, B] = 0$.

		\textit{Case 1.} Let $s \in \lbrace i, j \rbrace$, $t \in \lbrace i, j \rbrace$ and $k \not\in \lbrace i, j \rbrace$. Since $[B, A + E_{ks}] = [B, E_{ks}] \neq 0$, $[A, E_{ks}] \neq 0$ and $[A + B, E_{ks}] \neq 0$, we have
		\begin{equation*}
			Z(C(A + B)) \cap Z(C(E_{ks})) = Z(C(B)) \cap Z(C(A + E_{ks})) = Z(C(A)) \cap Z(C(E_{ks})) = \zzz.
		\end{equation*}
		Hence,
		\begin{equation*}
			\Delta(A, B) = \Delta(B, A + E_{ks}) + \Delta(A, E_{ks}) - \Delta(A + B, E_{ks}) \in \alpha(\zzz).
		\end{equation*}

		\textit{Case 2.} Let $s \neq i$ and $s \neq j$. Since $[B, A + E_{is}] = [B, E_{is}] \neq 0$, $[A, E_{is}] \neq 0$ and $[A + B, E_{is}] \neq 0$, we have
		\begin{equation*}
			Z(C(A + B)) \cap Z(C(E_{is})) = Z(C(B)) \cap Z(C(A + E_{is})) = Z(C(A)) \cap Z(C(E_{is})) = \zzz.
		\end{equation*}
		Hence,
		\begin{equation*}
			\Delta(A, B) = \Delta(B, A + E_{is}) + \Delta(A, E_{is}) - \Delta(A + B, E_{is}) \in \alpha(\zzz).
		\end{equation*}
		
		\textit{Case 3.} Let $t \neq i$ and $t \neq j$. Since $[B, A + E_{ti}] = [B, E_{ti}] \neq 0$, $[A, E_{ti}] \neq 0$ and $[A + B, E_{ti}] \neq 0$, we have
		\begin{equation*}
			Z(C(A + B)) \cap Z(C(E_{ti})) = Z(C(B)) \cap Z(C(A + E_{ti})) = Z(C(A)) \cap Z(C(E_{ti})) = \zzz.
		\end{equation*}
		Hence,
		\begin{equation*}
			\Delta(A, B) = \Delta(B, A + E_{ti}) + \Delta(A, E_{ti}) - \Delta(A + B, E_{ti}) \in \alpha(\zzz).
		\end{equation*}
	\end{proof}

\end{document}
